%
% pi2.tex -- template for standalon tikz images
%
% (c) 2021 Prof Dr Andreas Müller, OST Ostschweizer Fachhochschule
%
\documentclass[tikz]{standalone}
\usepackage{amsmath}
\usepackage{times}
\usepackage{txfonts}
\usepackage{pgfplots}
\usepackage{csvsimple}
\usetikzlibrary{arrows,intersections,math}
\begin{document}
\def\ys{0.5}
\def\skala{0.933}
\pgfmathparse{3.1415*3.1415}
\xdef\piq{\pgfmathresult}
\begin{tikzpicture}[>=latex,thick,scale=\skala,
declare function = {
f(\X) = 3.2899 + 4*(-1)*cos(deg(\X)) + 4*(0.25)*cos(deg(2*\X)) + 4*(-0.1111)*cos(deg(3*\X));
}]

\draw[color=gray,dashed] (3.1415,0) -- ++(0,{\ys*\piq});
\draw[dashed,color=gray] (0,{\ys*\piq}) -- ++(3.1415,0);
\foreach \y in {2,4,6,8}{
	\draw (-0.05,{\ys*\y}) -- ++(0.1,0);
	\node at (0,{\ys*\y}) [left] {$\y$};
}
\draw (-0.05,{\ys*\piq}) -- ++(0.1,0);
\node at (0,{\ys*\piq}) [left] {$\pi^2$};

\begin{scope}
	\clip (-6.28318,-0.2) rectangle (6.28318,{\ys*\piq});
	\draw[color=gray!50,line width=1.2pt]
		plot[domain=-3.1415:3.1415,samples=300] ({\x},{\ys*\x*\x)});

	\begin{scope}[xshift=6.28318cm]
		\draw[color=gray!50,line width=1.2pt]
			plot[domain=-3.1415:3.1415,samples=300]
				({\x},{\ys*\x*\x)});
	\end{scope}

	\begin{scope}[xshift=-6.28318cm]
		\draw[color=gray!50,line width=1.2pt]
			plot[domain=-3.1415:3.1415,samples=300]
				({\x},{\ys*\x*\x)});
	\end{scope}
\end{scope}
\node[color=gray!50] at (-3.1415,{\ys*\piq+0.1}) [right] {$y=x^2$};
\draw (3.1415,{\ys*\piq}) circle[radius=0.1];

\draw[->] (-6.4,0) -- (6.6,0) coordinate[label={$x$}];
\draw[->] (0,-0.5) -- (0,{\ys*11.0}) coordinate[label={right:$y$}];

\draw[color=blue!50,line width=1.2pt]
	plot[domain=-6.3:6.3,samples=300] ({\x},{\ys*f(\x)});
\node[color=blue!50] at (3.9,{\ys*(6.28-3.9)*(6.28-3.9)})
	[above right] {$y=f_3(x)$};

\node at (3.1415,-0.05) [below] {$\pi\mathstrut$};
\node at ({2*3.1415},-0.05) [below] {$2\pi\mathstrut$};
\draw (3.1415,-0.05) -- ++(0,0.1);
\draw ({2*3.1415},-0.05) -- ++(0,0.1);
\node at (-3.1415,-0.05) [below] {$-\pi\mathstrut$};
\node at ({-2*3.1415},-0.05) [below] {$-2\pi\mathstrut$};
\draw (-3.1415,-0.05) -- ++(0,0.1);
\draw ({-2*3.1415},-0.05) -- ++(0,0.1);

\end{tikzpicture}
\end{document}

